\section{Extracción de datos}
El componente de extracción consiste en un módulo independiente de Python encargado de la captura, procesamiento y 
empaquetado de la información interna del modelo. Este componente se articula en tres módulos funcionales: el 
\textbf{módulo de ejecución}, el \textbf{módulo de extracción} y el \textbf{módulo de formato e intercambio}.
\medskip

El \textbf{módulo de ejecución} contempla la arquitectura estándar para la carga y evaluación de la red neuronal 
convolucional. Dado que el alcance del proyecto se enfoca en la fase de inferencia, este módulo gestiona la 
inicialización de la topología y la carga de los pesos sinápticos pre-entrenados del modelo YOLOv2. El módulo 
recibe como entradas los datos de prueba (imágenes), los parámetros operativos y los pesos guardados, generando 
como resultado la predicción o detección del modelo.
\medskip

El \textbf{módulo de extracción} se encarga de interceptar y almacenar la información interna de la red durante 
el paso hacia adelante (\textit{forward pass}). Para ello, implementa ganchos de ejecución posterior 
(\textit{forward hooks}) en PyTorch, permitiendo capturar las representaciones intermedias, los mapas de activación 
y la topología de las capas, así como los metadatos contextuales de entrada y salida.
\medskip

Finalmente, el \textbf{módulo de formato e intercambio} transforma las estructuras matriciales a los formatos 
intermedios estandarizados. Este módulo organiza la topología y metadatos en un esquema estructurado para su 
serialización en un manifiesto \textit{JSON}, mientras que convierte los tensores de activación en arreglos 
multidimensionales compatibles con la librería \textit{pyopenvdb} para su exportación volumétrica. Asimismo, 
permite definir y el directorio de destino para los archivos generados.
\medskip

En síntesis, este componente recibe como insumos el modelo convolucional y los datos de entrada, produciendo 
como salida una estructura jerárquica de directorios que contiene el manifiesto de metadatos y las secuencias 
volumétricas de datos.